\documentclass[a4paper,11pt]{article}

\usepackage[margin=2.5cm]{geometry}
\usepackage{listings}
\usepackage{xcolor}
\usepackage{hyperref}
\usepackage{booktabs}
\usepackage{fancyvrb}

% ── Colour palette ──
\definecolor{kw}{HTML}{A6261D}
\definecolor{str}{HTML}{6E8B3D}
\definecolor{cmt}{HTML}{888888}
\definecolor{repl}{HTML}{2E8B57}
\definecolor{bg}{HTML}{F7F7F7}
\definecolor{link}{HTML}{2A5DB0}
\definecolor{shell}{HTML}{4A148C}
\definecolor{nix-kw}{HTML}{7C4DFF}

\hypersetup{colorlinks=true,linkcolor=link,urlcolor=link}

% ── Listings: Common Lisp ──
\lstdefinelanguage{CL}{
  morekeywords={defun,defvar,defparameter,defmacro,lambda,let,let*,%
    if,cond,when,unless,progn,loop,do,dolist,dotimes,%
    setf,setq,funcall,apply,mapcar,format,print,load,%
    in-package,defpackage,export,use-package,%
    t,nil,quote},
  sensitive=true,
  morestring=[b]",
  morecomment=[l]{;},
  morecomment=[n]{\#|}{|\#},
}

\lstset{
  language=CL,
  basicstyle=\ttfamily\small,
  keywordstyle=\color{kw}\bfseries,
  stringstyle=\color{str},
  commentstyle=\color{cmt}\itshape,
  backgroundcolor=\color{bg},
  frame=single,
  rulecolor=\color{black!20},
  framesep=6pt,
  xleftmargin=12pt,
  xrightmargin=12pt,
  breaklines=true,
  showstringspaces=false,
  tabsize=2,
  columns=flexible,
}

% ── Shell style ──
\lstdefinestyle{shell}{
  language={},
  basicstyle=\ttfamily\small\color{shell},
  backgroundcolor=\color{bg},
  frame=single,
  rulecolor=\color{black!20},
  framesep=6pt,
  xleftmargin=12pt,
  xrightmargin=12pt,
  showstringspaces=false,
}

% ── REPL style ──
\lstdefinestyle{repl}{
  language=CL,
  basicstyle=\ttfamily\small,
  backgroundcolor=\color{bg},
  frame=single,
  rulecolor=\color{black!20},
  framesep=6pt,
  xleftmargin=12pt,
  xrightmargin=12pt,
  keywordstyle=\color{kw}\bfseries,
  stringstyle=\color{str},
  commentstyle=\color{cmt}\itshape,
  moredelim=[is][\color{repl}\bfseries]{|r|}{|r|},
  showstringspaces=false,
}

% ── Listings: Nix ──
\lstdefinelanguage{Nix}{
  morekeywords={let,in,if,then,else,with,rec,inherit,import,true,false,null},
  sensitive=true,
  morestring=[b]",
  morestring=[b]'',
  morecomment=[l]{\#},
  morecomment=[n]{/*}{*/},
}

\lstdefinestyle{nix}{
  language=Nix,
  basicstyle=\ttfamily\small,
  keywordstyle=\color{nix-kw}\bfseries,
  stringstyle=\color{str},
  commentstyle=\color{cmt}\itshape,
  backgroundcolor=\color{bg},
  frame=single,
  rulecolor=\color{black!20},
  framesep=6pt,
  xleftmargin=12pt,
  xrightmargin=12pt,
  breaklines=true,
  showstringspaces=false,
  tabsize=2,
  columns=flexible,
}

\newcommand{\key}[1]{\texttt{#1}}
\newcommand{\file}[1]{\texttt{#1}}
\newcommand{\pkg}[1]{\texttt{#1}}

% ── Side-by-side layout helpers ──
\newlength{\colwidth}
\setlength{\colwidth}{0.47\textwidth}

\newcommand{\colheader}[1]{\textbf{\large #1}\par\medskip}

% ── Narrow listing style (for minipage columns) ──
\lstdefinestyle{narrow}{
  basicstyle=\ttfamily\footnotesize,
  xleftmargin=4pt,
  xrightmargin=4pt,
  framesep=4pt,
  breaklines=true,
}

\title{\textbf{Common Lisp Workflow Tutorial}\\[4pt]
  \large REPL, Running Files, and Editor Integration}
\author{CS\,3339 --- Computer Architecture}
\date{\today}

\begin{document}
\maketitle
\tableofcontents
\newpage

%% ════════════════════════════════════════════════════════════════
\section{Setup}
%% ════════════════════════════════════════════════════════════════

\subsection{Install SBCL}

\begin{lstlisting}[style=shell]
# macOS
brew install sbcl

# Ubuntu / Debian
sudo apt install sbcl

# Nix
nix-env -iA nixpkgs.sbcl
\end{lstlisting}

Verify:

\begin{lstlisting}[style=shell]
$ sbcl --version
SBCL 2.4.x
\end{lstlisting}

%% ════════════════════════════════════════════════════════════════
\section{Using the REPL (Interactive Execution)}
%% ════════════════════════════════════════════════════════════════

REPL = Read--Eval--Print Loop.
You type an expression, it evaluates immediately, and prints the result.

\subsection{Step 1: Launch the REPL}

\begin{lstlisting}[style=shell]
$ sbcl
\end{lstlisting}

The prompt changes to \texttt{*}. This is the SBCL REPL.

\subsection{Step 2: Type expressions and run them}

\begin{lstlisting}[style=repl]
* (+ 1 2)
|r|3|r|

* (format t "Hello, world!~%")
|r|Hello, world!
NIL|r|

* (defun square (x) (* x x))
|r|SQUARE|r|

* (square 7)
|r|49|r|
\end{lstlisting}

Each expression is evaluated immediately and the result is printed.

\subsection{Step 3: Quit the REPL}

\begin{lstlisting}[style=repl]
* (quit)
\end{lstlisting}

Or press \key{Ctrl+D}.

%% ════════════════════════════════════════════════════════════════
\section{Writing Code in a File and Running It}
%% ════════════════════════════════════════════════════════════════

\subsection{Step 1: Create a \file{.lisp} file}

Open your editor and create \file{hello.lisp}:

\begin{lstlisting}
;;; hello.lisp --- first program

(defun greet (name)
  (format t "Hello, ~A!~%" name))

(greet "World")
\end{lstlisting}

\subsection{Step 2: Run the file from the terminal}

\textbf{Method A} --- \texttt{sbcl --script} (simplest, just runs and exits):

\begin{lstlisting}[style=shell]
$ sbcl --script hello.lisp
Hello, World!
\end{lstlisting}

\textbf{Method B} --- \texttt{sbcl --load} (runs file, then drops into REPL):

\begin{lstlisting}[style=shell]
$ sbcl --load hello.lisp
Hello, World!
* (greet "Alice")
Hello, Alice!
\end{lstlisting}

With \texttt{--load}, you stay in the REPL after execution,
so you can interactively test the functions you defined.

\textbf{Method C} --- \texttt{load} from inside the REPL:

\begin{lstlisting}[style=repl]
$ sbcl
* (load "hello.lisp")
|r|Hello, World!
T|r|
* (greet "Bob")
|r|Hello, Bob!
NIL|r|
\end{lstlisting}

\subsection{Step 3: Edit and reload}

After modifying \file{hello.lisp}, just \texttt{load} it again in the REPL:

\begin{lstlisting}[style=repl]
* (load "hello.lisp")   ; changes take effect immediately
\end{lstlisting}

\subsection{Summary of methods}

\begin{center}
\begin{tabular}{lll}
  \toprule
  \textbf{Method} & \textbf{Command} & \textbf{Use case} \\
  \midrule
  Script        & \texttt{sbcl --script file.lisp} & One-shot execution \\
  Load + REPL   & \texttt{sbcl --load file.lisp}   & Run then interact \\
  REPL load     & \texttt{(load "file.lisp")}      & Reload during dev \\
  \bottomrule
\end{tabular}
\end{center}

%% ════════════════════════════════════════════════════════════════
\section{REPL-Driven Development}
%% ════════════════════════════════════════════════════════════════

The Common Lisp development cycle is very tight:
write code, send it to the REPL, test it, fix it --- all without restarting.

\begin{enumerate}
  \item Write code in a \file{.lisp} file in your editor.
  \item Send the file (or individual functions) to the REPL.
  \item Test interactively in the REPL.
  \item Go back to the editor, fix things, and repeat from step 2.
\end{enumerate}

Editors like SLIME (Emacs) and Vlime (Neovim) let you do steps 1--3
without ever leaving your editor. That is their main advantage.

%% ════════════════════════════════════════════════════════════════
\section{Editor Integration}
%% ════════════════════════════════════════════════════════════════

Both Emacs (with SLIME) and Neovim (with Vlime) connect to a running SBCL
and let you evaluate code, look up documentation, and debug --- all without
leaving the editor. The subsections below compare the two side by side.

%% ── 5.1 ──
\subsection{Installation}

\noindent
\begin{minipage}[t]{\colwidth}
  \colheader{Emacs + SLIME}
  Add to \file{\textasciitilde/.emacs.d/init.el}:
\begin{lstlisting}[style=shell,style=narrow]
(require 'package)
(add-to-list 'package-archives
 '("melpa" .
   "https://melpa.org/packages/"))
(package-initialize)
;; M-x package-install RET slime RET

(setq inferior-lisp-program "sbcl")
\end{lstlisting}
\end{minipage}\hfill
\begin{minipage}[t]{\colwidth}
  \colheader{Neovim + Vlime}
  Add to \file{\textasciitilde/.config/nvim/init.vim}:
\begin{lstlisting}[style=shell,style=narrow]
Plug 'vlime/vlime', {'rtp': 'vim/'}
\end{lstlisting}
  Then run \texttt{:PlugInstall} inside Neovim.
\end{minipage}

%% ── 5.2 ──
\subsection{Starting the REPL Connection}

\noindent
\begin{minipage}[t]{\colwidth}
  \colheader{Emacs}
\begin{lstlisting}[style=shell,style=narrow]
M-x slime
\end{lstlisting}
  A \texttt{CL-USER>} prompt appears at the bottom of the screen.
\end{minipage}\hfill
\begin{minipage}[t]{\colwidth}
  \colheader{Neovim}
  Press \key{\textbackslash rr} inside Neovim to start SBCL and open a REPL buffer.

  \medskip
  Alternatively, start the server externally:
\begin{lstlisting}[style=shell,style=narrow]
$ vlime-server
\end{lstlisting}
  Then connect from Neovim with \key{\textbackslash rc}.
\end{minipage}

%% ── 5.3 ──
\subsection{Vlime Window Layout}

After pressing \key{\textbackslash rr}, Neovim splits into several buffers.
Each serves a distinct role:

\begin{center}
\begin{tabular}{llp{7.5cm}}
  \toprule
  \textbf{Status bar} & \textbf{Name} & \textbf{Role} \\
  \midrule
  \texttt{vlime | repl}   & REPL buffer
    & Displays compilation results (e.g.\ \texttt{FIB}) and
      return values of evaluated expressions. \\
  \texttt{vlime | server} & Server buffer
    & Shows raw SBCL and Swank startup messages.
      Useful for diagnosing connection issues. \\
  \texttt{vlime | sldb}   & Debugger (SLDB)
    & Appears when an error occurs.
      Lists restarts and a stack-frame backtrace.
      Press \key{\textbackslash a} to abort and close. \\
  \texttt{vlime | input}  & Input buffer
    & Opens when using \key{\textbackslash ss} etc.
      Edit the expression, then press \key{Enter} to send it to the REPL. \\
  \bottomrule
\end{tabular}
\end{center}

Your source file (e.g.\ \file{tutorial.lisp}) remains in its own buffer
alongside these Vlime windows.

In Emacs, SLIME uses a similar split: the \texttt{*slime-repl*} buffer for the
REPL and the \texttt{*sldb*} buffer for the debugger.

%% ── 5.4 ──
\subsection{Evaluating Code}

\begin{center}
\begin{tabular}{lll}
  \toprule
  \textbf{Action} & \textbf{Emacs (SLIME)} & \textbf{Neovim (Vlime)} \\
  \midrule
  Evaluate expression      & \key{C-x C-e}  & \key{\textbackslash ss} \\
  Evaluate top-level form   & \key{C-M-x}    & \key{\textbackslash se} \\
  Evaluate region/selection & \key{C-c C-r}  & \key{\textbackslash sr} \\
  Compile defun             & \key{C-c C-c}  & --- \\
  Compile \& load file      & \key{C-c C-l}  & \key{\textbackslash ol} \\
  Switch to REPL            & \key{C-c C-z}  & (REPL buffer) \\
  \bottomrule
\end{tabular}
\end{center}

\paragraph{Vlime \key{\textbackslash ss} workflow.}
Pressing \key{\textbackslash ss} opens an input buffer (\texttt{vlime | input}).
The expression under the cursor is pre-filled.
Edit it if needed, then press \key{Enter} to send it to the REPL.
The result appears in the REPL buffer.
\key{\textbackslash se} skips the input buffer and evaluates the top-level
form directly.

\paragraph{Note on the Vlime REPL buffer.}
Unlike Emacs SLIME, where you can type directly into the \texttt{*slime-repl*}
buffer, the Vlime REPL buffer is \emph{read-only} --- it only displays output.
To send an arbitrary expression to the REPL, use \key{\textbackslash si}
(interactive input), which opens an input buffer for you to type in.

%% ── 5.5 ──
\subsection{Documentation and Debugging}

\begin{center}
\begin{tabular}{lll}
  \toprule
  \textbf{Action} & \textbf{Emacs (SLIME)} & \textbf{Neovim (Vlime)} \\
  \midrule
  Describe symbol           & \key{C-c C-d d} & \key{\textbackslash dd} \\
  Open HyperSpec            & \key{C-c C-d h} & \key{\textbackslash dh} \\
  Show arglist              & \key{C-c C-d a} & \key{\textbackslash do} \\
  Macroexpand               & \key{C-c C-m}   & \key{\textbackslash me} \\
  Trace function            & \key{C-c C-t}   & \key{\textbackslash tf} \\
  \bottomrule
\end{tabular}
\end{center}

%% ── 5.5 ──
\subsection{Walkthrough Example}

\begin{enumerate}
  \item Open \file{fib.lisp}:

\begin{lstlisting}
(defun fib (n)
  "Return the nth Fibonacci number."
  (if (< n 2)
      n
      (+ (fib (- n 1))
         (fib (- n 2)))))
\end{lstlisting}

  \item Place cursor on the \texttt{defun} and evaluate it
        (\key{C-c C-c} in Emacs, \key{\textbackslash se} in Vlime).
        The REPL shows \texttt{FIB} --- the function is now defined.

  \item Switch to the REPL
        (\key{C-c C-z} in Emacs, or move to the REPL buffer in Neovim)
        and test:

\begin{lstlisting}[style=repl]
CL-USER> (fib 10)
|r|55|r|
CL-USER> (fib 30)
|r|832040|r|
\end{lstlisting}

  \item Go back to the file, edit, and re-evaluate.
        No recompilation or restart needed --- changes apply instantly.
\end{enumerate}

%% ════════════════════════════════════════════════════════════════
\section{Full Walkthrough: From File to Execution}
%% ════════════════════════════════════════════════════════════════

A minimal example using only the terminal (no editor plugins needed).

\subsection{Step 1: Create a project directory}

\begin{lstlisting}[style=shell]
$ mkdir ~/cl-practice && cd ~/cl-practice
\end{lstlisting}

\subsection{Step 2: Write the program}

Create \file{fizzbuzz.lisp}:

\begin{lstlisting}
;;; fizzbuzz.lisp

(defun fizzbuzz (n)
  (loop for i from 1 to n
        do (format t "~A~%"
             (cond
               ((zerop (mod i 15)) "FizzBuzz")
               ((zerop (mod i 3))  "Fizz")
               ((zerop (mod i 5))  "Buzz")
               (t                   i)))))

;; This runs when the file is loaded
(fizzbuzz 20)
\end{lstlisting}

\subsection{Step 3: Run it}

\begin{lstlisting}[style=shell]
$ sbcl --script fizzbuzz.lisp
1
2
Fizz
4
Buzz
Fizz
7
8
Fizz
Buzz
11
Fizz
13
14
FizzBuzz
16
17
Fizz
19
Buzz
\end{lstlisting}

\subsection{Step 4: Interactive testing}

\begin{lstlisting}[style=shell]
$ sbcl --load fizzbuzz.lisp
\end{lstlisting}

After the output, the REPL is ready:

\begin{lstlisting}[style=repl]
* (fizzbuzz 5)
|r|1
2
Fizz
4
Buzz
NIL|r|
\end{lstlisting}

Change the argument, call it again --- no restart required.

%% ════════════════════════════════════════════════════════════════
\section{Quick Reference}
%% ════════════════════════════════════════════════════════════════

\begin{center}
\begin{tabular}{lll}
  \toprule
  \textbf{Action} & \textbf{Emacs (SLIME)} & \textbf{Neovim (Vlime)} \\
  \midrule
  \multicolumn{3}{l}{\textbf{--- Terminal (both editors) ---}} \\
  Start the REPL             & \multicolumn{2}{c}{\texttt{sbcl}} \\
  Run a file as a script     & \multicolumn{2}{c}{\texttt{sbcl --script file.lisp}} \\
  Load a file then REPL      & \multicolumn{2}{c}{\texttt{sbcl --load file.lisp}} \\
  Reload inside REPL         & \multicolumn{2}{c}{\texttt{(load "file.lisp")}} \\
  Quit the REPL              & \multicolumn{2}{c}{\texttt{(quit)} or \key{Ctrl+D}} \\
  \midrule
  \multicolumn{3}{l}{\textbf{--- Editor ---}} \\
  Start REPL connection      & \key{M-x slime}  & \key{\textbackslash rr} \\
  Evaluate expression        & \key{C-x C-e}    & \key{\textbackslash ss} \\
  Evaluate top-level form    & \key{C-M-x}      & \key{\textbackslash se} \\
  Compile defun              & \key{C-c C-c}    & --- \\
  Compile \& load file       & \key{C-c C-l}    & \key{\textbackslash ol} \\
  Switch to REPL             & \key{C-c C-z}    & (REPL buffer) \\
  Describe symbol            & \key{C-c C-d d}  & \key{\textbackslash dd} \\
  Macroexpand                & \key{C-c C-m}    & \key{\textbackslash me} \\
  \bottomrule
\end{tabular}
\end{center}

\newpage
\appendix

%% ════════════════════════════════════════════════════════════════
\section{flake.nix --- Development Environment}
%% ════════════════════════════════════════════════════════════════

This project uses \textbf{Nix Flakes} to provide a reproducible development environment.
The \file{flake.nix} file declaratively specifies all dependencies needed for the project:
a Common Lisp compiler (SBCL) with libraries, a Vlime server for Neovim integration,
and a \LaTeX{} distribution for document compilation.

Running \key{nix develop} in the project root activates a shell
with all tools available, without installing anything globally.

\subsection{Prerequisites}

\begin{enumerate}
  \item \textbf{Nix} (version 2.4 or later) with the \texttt{flakes} and
        \texttt{nix-command} experimental features enabled.
  \item Add the following to \file{\textasciitilde/.config/nix/nix.conf}
        (or \file{/etc/nix/nix.conf}):
\begin{lstlisting}[style=shell]
experimental-features = nix-command flakes
\end{lstlisting}
\end{enumerate}

\subsection{Usage}

\subsubsection{Entering the development shell}

\begin{lstlisting}[style=shell]
$ cd /path/to/project
$ nix develop
\end{lstlisting}

This drops you into a shell where \pkg{sbcl}, \pkg{vlime-server},
and \pkg{pdflatex} are all available.

\subsubsection{Running Common Lisp code}

\begin{lstlisting}[style=shell]
$ sbcl --script tutorial.lisp
$ sbcl --load tutorial.lisp
\end{lstlisting}

\subsubsection{Starting the Vlime server}

\begin{lstlisting}[style=shell]
$ vlime-server
\end{lstlisting}

Then connect from Neovim with \key{\textbackslash cc}.

\subsubsection{Compiling \LaTeX{} documents}

\begin{lstlisting}[style=shell]
$ pdflatex doc/tutorial.tex
\end{lstlisting}

\subsection{Structure of \file{flake.nix}}

The flake is composed of two top-level attributes: \texttt{inputs} and \texttt{outputs}.

\subsubsection{Inputs}

\begin{lstlisting}[style=nix]
inputs = {
  nixpkgs.url = "github:nixos/nixpkgs/nixpkgs-unstable";
};
\end{lstlisting}

The sole input is \textbf{nixpkgs} from the \texttt{nixpkgs-unstable} branch.
This provides access to the latest available packages.
The exact revision is pinned in \file{flake.lock} to ensure reproducibility.

\subsubsection{Outputs}

The \texttt{outputs} function receives \texttt{self} (the flake itself) and
\texttt{nixpkgs} (the resolved input), then produces a set of \texttt{devShells}.

\subsubsection{Multi-platform support}

\begin{lstlisting}[style=nix]
forAllSystems = nixpkgs.lib.genAttrs [
  "aarch64-darwin"
  "x86_64-darwin"
  "x86_64-linux"
  "aarch64-linux"
];
\end{lstlisting}

The helper \texttt{forAllSystems} generates outputs for four platforms:

\begin{center}
\begin{tabular}{ll}
  \toprule
  \textbf{System} & \textbf{Platform} \\
  \midrule
  \texttt{aarch64-darwin} & macOS (Apple Silicon) \\
  \texttt{x86\_64-darwin} & macOS (Intel) \\
  \texttt{x86\_64-linux}  & Linux (Intel/AMD) \\
  \texttt{aarch64-linux}  & Linux (ARM) \\
  \bottomrule
\end{tabular}
\end{center}

\subsection{Packages Provided}

The development shell provides three components.

\subsubsection{SBCL with Common Lisp libraries}

\begin{lstlisting}[style=nix]
sbclWithDeps = pkgs.sbcl.withPackages (ps: [
  ps.swank
  ps.alexandria
  ps.yason
  ps.vom
]);
\end{lstlisting}

\textbf{SBCL} (Steel Bank Common Lisp) is bundled with the following libraries:

\begin{center}
\begin{tabular}{lp{9cm}}
  \toprule
  \textbf{Library} & \textbf{Description} \\
  \midrule
  \pkg{swank}      & Backend server for SLIME/Vlime;
                     enables REPL-driven development from an editor. \\
  \pkg{alexandria}  & A widely-used utility library providing common functions
                     (e.g.\ hash-table operations, list utilities, type checks). \\
  \pkg{yason}      & A JSON encoder/decoder for Common Lisp. \\
  \pkg{vom}        & A logging library with log-level support. \\
  \bottomrule
\end{tabular}
\end{center}

\subsubsection{Vlime server}

\begin{lstlisting}[style=nix]
vlime-src = pkgs.fetchFromGitHub {
  owner = "vlime";
  repo = "vlime";
  rev = "e276e9a6f37d2699a3caa63be19314f5a19a1481";
  hash = "sha256-tCqN80lgj11ggzGmuGF077oqL5ByjUp6jVmRUTrIWJA=";
};

vlime-server = pkgs.writeShellScriptBin "vlime-server" ''
  exec ''${sbclWithDeps}/bin/sbcl --noinform \
    --load ''${vlime-src}/lisp/start-vlime.lisp
'';
\end{lstlisting}

A wrapper script \pkg{vlime-server} is created that:

\begin{enumerate}
  \item Fetches the Vlime source from GitHub at a pinned commit.
  \item Launches SBCL (with all libraries loaded) and starts the Vlime server.
  \item Suppresses the SBCL startup banner with \texttt{--noinform}.
\end{enumerate}

This provides a single command to start the Vlime server for Neovim integration,
without requiring any manual installation of Vlime.

\subsubsection{\LaTeX{} distribution}

\begin{lstlisting}[style=nix]
pkgs.texliveMedium
\end{lstlisting}

\pkg{texliveMedium} is a mid-size TeX Live distribution that includes
common packages such as \pkg{listings}, \pkg{xcolor}, \pkg{hyperref},
\pkg{booktabs}, and \pkg{geometry} --- everything needed to compile
the project's documentation.

\subsection{The \file{flake.lock} File}

The \file{flake.lock} file is automatically generated and records the exact
revision of each input. This ensures that every user gets the same package
versions regardless of when they run \key{nix develop}.

To update all inputs to their latest versions:

\begin{lstlisting}[style=shell]
$ nix flake update
\end{lstlisting}

To update only nixpkgs:

\begin{lstlisting}[style=shell]
$ nix flake update nixpkgs
\end{lstlisting}

After updating, commit both \file{flake.nix} and \file{flake.lock}.

\subsection{Quick Reference}

\begin{center}
\begin{tabular}{lp{8cm}}
  \toprule
  \textbf{Command} & \textbf{Description} \\
  \midrule
  \texttt{nix develop}            & Enter the development shell \\
  \texttt{sbcl}                   & Start the Common Lisp REPL \\
  \texttt{sbcl --script file.lisp}& Run a Lisp file as a script \\
  \texttt{vlime-server}           & Start the Vlime server for Neovim \\
  \texttt{pdflatex doc/file.tex}  & Compile a \LaTeX{} document \\
  \texttt{nix flake update}       & Update all inputs to latest \\
  \texttt{nix flake show}         & Show flake outputs \\
  \texttt{nix flake metadata}     & Show flake metadata and lock info \\
  \bottomrule
\end{tabular}
\end{center}

\end{document}
